%% Chapter 6 — Evaluation
\chapter{Evaluation}
\label{ch:evaluation}

This chapter evaluates the implemented NPR avatar system and documents the planned user study for trust assessment.

\section{Technical Evaluation Setup}
\label{sec:technical_setup}

\subsection{Evaluation Objectives}

The technical evaluation verifies:
\begin{itemize}
  \item the quality and coherence of NPR avatar rendering in VR
  \item real-time performance and stability under realistic usage
  \item readiness of the system for a trust-focused user study.
\end{itemize}

\subsection{Technical Metrics}

Define the quantitative and qualitative metrics:
\begin{itemize}
  \item average frame rate and minimum frame rate
  \item frame time variance and rendering latency
  \item stylization fidelity indicators such as edge coherence and temporal stability
  \item subjective observations from developer-led test sessions.
\end{itemize}

\subsection{Baseline Comparisons}

List the technical baselines used for comparison:
\begin{itemize}
  \item non-stylized VR avatar rendering
  \item existing NPR avatar rendering approaches, if available
  \item simplified style variants to isolate the effect of each component.
\end{itemize}

\section{Implementation Validation}
\label{sec:implementation_validation}

\subsection{Performance Results}

Summarize measured performance:
\begin{itemize}
  \item frame rate and latency results in representative VR scenarios
  \item resource usage on target hardware
  \item performance behavior for each style variant.
\end{itemize}

\subsection{Visual Quality Assessment}

Report qualitative observations and comparisons:
\begin{itemize}
  \item visual consistency across head motion
  \item edge quality, shading coherence, and stylization clarity
  \item artifacts or failure cases observed during testing.
\end{itemize}

\subsection{Stability and Robustness}

Describe the stability of the implementation:
\begin{itemize}
  \item behavior during extended VR sessions
  \item integration issues and debugging outcomes
  \item support for varying avatar animations and scene complexity.
\end{itemize}

\section{User Study Design}
\label{sec:user_study_design}

This section describes the planned user study that will assess how NPR avatar styles influence trust.

\subsection{Study Goals}

The study aims to answer:
\begin{itemize}
  \item how do different NPR avatar styles affect perceived trustworthiness in VR?
  \item which visual traits most strongly influence trust and social presence?
  \item does performance or visual quality moderate trust judgments?
\end{itemize}

\subsection{Participants and Recruitment}

Define the participant pool and recruitment plan:
\begin{itemize}
  \item target sample size and demographics
  \item inclusion and exclusion criteria
  \item ethical considerations and informed consent.
\end{itemize}

\subsection{Experimental Conditions}

Describe the conditions participants will experience:
\begin{itemize}
  \item baseline avatar rendering
  \item multiple NPR avatar variants with different stylization levels
  \item within-subject or between-subject study design.
\end{itemize}

\subsection{Task Scenario}

Outline the interaction scenario used in the study:
\begin{itemize}
  \item task or dialog with the avatar
  \item session duration and condition order
  \item behavioural measures such as task completion or response time.
\end{itemize}

\subsection{Measurement Instruments}

List the trust and presence instruments:
\begin{itemize}
  \item trust questionnaire adapted from VR or HRI research
  \item presence / immersion scale
  \item post-condition ratings for comfort, realism, and trust
  \item optional behavioural metrics such as interaction choices.
\end{itemize}

\subsection{Pilot Study Plan}

Describe the pilot study to validate the protocol:
\begin{itemize}
  \item a small pilot participant set
  \item testing condition order, questionnaire clarity, and timing
  \item refining instructions and VR content before the full study.
\end{itemize}

\section{Discussion}
\label{sec:discussion}

Reflect on the implementation and readiness for the user study.

\subsection{Readiness for the User Study}

Summarize whether the current system is ready for the planned study and what remains to be done.

\subsection{Limitations}

Discuss limitations affecting the planned study:
\begin{itemize}
  \item hardware and platform constraints
  \item asset and style coverage limitations
  \item potential bias in the selected stylization conditions.
\end{itemize}
